\section{System Model}
\label{sec:model}

A guest bursts batch GPU Jobs onto a shared cluster and chooses a self-imposed
concurrency window $K$. Time is slotted into control epochs $t\in\mathbb{N}$. Let
$a_t\in\{0,\dots,C\}$ be the concurrency that can produce useful guest work when
the action for epoch $t$ takes effect. The process is set by other users and the
scheduler, so the model treats it as exogenous and not directly observable. The
guest chooses an executed window $w_t\in\{0,\dots,K\}$. Useful goodput is
$G_t=\min(w_t,a_t)$, and $(w_t-a_t)^+$ is capacity mismatch. Depending on the
scheduler, excess work can remain pending or can overlap with capacity that is
not productive for the guest. The model does not equate every mismatch with an
administrative policy violation.

Writing
$\rho\triangleq\limsup_T T^{-1}\sum_{t<T}\mathbf 1[w_t>a_t]$ for the long-run
mismatch fraction, the guest solves
\begin{equation}\label{eq:polite}
\begin{aligned}
\max_{\{w_t\}}\quad & \mathbb{E}[\min(w_t,a_t)] \\
\text{s.t.}\quad & w_t\le K\ \forall t,\qquad \rho\le\epsilon .
\end{aligned}
\end{equation}
The hard window and the mismatch budget represent different constraints. The
first holds pathwise. The second compares policies at a common capacity-risk
level.

\noindent\textbf{The observation-to-effect lag.}
A decision at wall-clock time $\tau$ uses a capacity record sampled
$D_{\mathrm{obs}}$ earlier. A launch-on-demand action becomes productive at
$\tau+D_{\mathrm{act}}$. The state that determines its payoff is therefore the
capacity at $\tau+D_{\mathrm{act}}$, while its information was sampled at
$\tau-D_{\mathrm{obs}}$. We define
\begin{equation}\label{eq:Ddef}
D\ \triangleq\ D_{\mathrm{obs}}+D_{\mathrm{act}}
\end{equation}
as the observation-to-effect lag, expressed in decision epochs. Relabelling the
effect time as $t$, a lag-$D$ policy is causal with respect to
$\{a_{s-D}\}_{s\le t}$. The capacity record at decision time is only
$D_{\mathrm{obs}}$ old. The additional $D_{\mathrm{act}}$ appears because the
admission is scored when its effect is realized.

The measurements in Section~\ref{sec:eval} separate these quantities. Polling
already running resources takes $27$ to $33$ ms at the median. The measured path
from launch to visible useful GPU work takes about $2.46$ s. For the
launch-on-demand implementation, actuation dominates the lag. This need not hold
for a warm pool or a persistent worker queue.

\begin{proposition}[Launch-on-demand effect bound]
\label{prop:launchbound}
Suppose a scale-out action is realized only when a worker created by that action
becomes ready, and let $\tau_{\mathrm{ready}}$ be that worker's launch-to-ready
time. Then the action path has
$D_{\mathrm{act}}\ge\tau_{\mathrm{ready}}$ and
$D\ge D_{\mathrm{obs}}+\tau_{\mathrm{ready}}$.
\end{proposition}
\begin{proof}
The action cannot produce work before the worker that realizes it is ready.
Substitution into \eqref{eq:Ddef} gives the second inequality.
\end{proof}
Proposition~\ref{prop:launchbound} is implementation scoped. It does not rule out
warm workers, pipelining or speculative reservations. Those mechanisms can trade
resident resources or extra offers for a smaller effect delay and would require
a model that accounts for that cost.

\noindent\textbf{The contended unit.}
Politeness is decided at the margin, so we first isolate one fluctuating resource
unit. Let $x_t\in\{0,1\}$ indicate that the unit is free at effect time. The
guest's marginal decision $m_t\in\{0,1\}$ is a possibly randomized causal
function of the lag-$D$ observation
$\hat x_t\triangleq x_{t-D}$. We take $x_t$ stationary with
$\Pr[x_t=1]=1/2$. Marginal goodput and mismatch are
\begin{equation}\label{eq:margin}
g \triangleq \Pr[m_t=1,\,x_t=1],\qquad
\rho \triangleq \Pr[m_t=1,\,x_t=0].
\end{equation}
The stationary $\rho$ in \eqref{eq:margin} is the ergodic counterpart of the
epoch fraction in \eqref{eq:polite}.

\begin{proposition}[Window safety]\label{prop:safety}
Any update that clamps the executed window to $[0,K]$ keeps $w_t\le K$ for every
$t$ deterministically.
\end{proposition}
\begin{proof}
The claim follows directly from the clamp at every update.
\end{proof}

\section{How Much Delayed Feedback Is Worth}
\label{sec:limit}

We characterize the gain over a record-independent policy when the only state
signal is a lag-$D$ observation.

\begin{lemma}[Predictive disagreement]\label{lem:q}
Let $q(D)\triangleq\Pr[x_t\neq x_{t-D}]$ and
$r(D)\triangleq\mathbb{E}[(2x_t-1)(2x_{t-D}-1)]$. Then
$r(D)=1-2q(D)$. Under stationarity and symmetry,
\[
\Pr[\hat x_t=i,x_t=j]
=\tfrac14\bigl(1+(2i-1)(2j-1)r(D)\bigr).
\]
For a symmetric two-state Markov process with per-epoch flip probability
$p\in(0,1/2)$,
$r(D)=(1-2p)^D=e^{-D/\tau_c}$, where
$\tau_c=-1/\ln(1-2p)$.
\end{lemma}
\begin{proof}
The correlation equals agreement probability minus disagreement probability.
The joint law follows from the two uniform marginals. For the Markov process, the
second eigenvalue is $1-2p$, so the lag-$D$ correlation is its $D$th power.
\end{proof}

\begin{theorem}[Matched-mismatch frontier]\label{thm:limit}
For the symmetric two-state Markov process of Lemma~\ref{lem:q}, the achievable
$(\rho,g)$ region under lag $D$ is the quadrilateral with vertices
\[
(0,0),\quad
\left(\tfrac q2,\tfrac{1-q}{2}\right),\quad
\left(\tfrac12,\tfrac12\right),\quad
\left(\tfrac{1-q}{2},\tfrac q2\right),
\]
where $q=q(D)$. Its efficient upper boundary joins the first three vertices. A
record-independent policy is confined to $g=\rho$. Consequently the maximum
vertical goodput increase at matched mismatch is
\begin{equation}\label{eq:adv}
\Delta(D)=\tfrac12 r(D)=\tfrac12\bigl(1-2q(D)\bigr).
\end{equation}
The maximum occurs at $\rho=q/2$ and is attained by the threshold policy
$m_t=\hat x_t$. For $0\le\rho<q/2$, the increase is
$\rho(1-2q)/q$.
\end{theorem}
\begin{proof}
The Markov property makes $\hat x_t$ sufficient for $x_t$ given the full delayed
history. Any causal policy induces
$\theta_i=\Pr[m_t=1\mid\hat x_t=i]$. Lemma~\ref{lem:q} gives
\[
g=\tfrac12\bigl((1-q)\theta_1+q\theta_0\bigr),\qquad
\rho=\tfrac12\bigl(q\theta_1+(1-q)\theta_0\bigr).
\]
The image of $[0,1]^2$ under this linear map is the stated quadrilateral. Since
$q<1/2$, admission after $\hat x_t=1$ has the larger goodput-to-mismatch ratio,
so the upper boundary first increases $\theta_1$ and then $\theta_0$. A
record-independent policy has $\theta_1=\theta_0$, hence $g=\rho$. The vertical
difference at the kink is $(1-2q)/2=r(D)/2$.
\end{proof}

\begin{figure}[t]
\centering
\resizebox{0.92\columnwidth}{!}{%
\begin{tikzpicture}[>=Stealth]
  \draw[->] (0,0) -- (6.0,0) node[right] {$\rho$ (mismatch)};
  \draw[->] (0,0) -- (0,6.0) node[above] {$g$ (goodput)};
  \fill[black!8] (0,0) -- (1,4) -- (5,5) -- (4,1) -- cycle;
  \node[black!55] at (4.05,1.55) {\footnotesize achievable};
  \draw[very thick] (0,0) -- (1,4) -- (5,5);
  \draw[dashed] (0,0) -- (5.4,5.4);
  \node[sloped,above,black!70] at (3.25,3.25) {\footnotesize record-independent, $g=\rho$};
  \draw[dotted] (1,0) -- (1,4);
  \node[below] at (1,0) {\footnotesize $\tfrac{q}{2}$};
  \fill (1,4) circle (2pt);
  \node[above left=-1pt] at (1,4) {\footnotesize $\bigl(\tfrac{q}{2},\tfrac{1-q}{2}\bigr)$};
  \node[align=center,anchor=south west] at (1.55,4.45)
        {\footnotesize delayed threshold\\[-2pt]\footnotesize $m_t=\hat x_t$};
  \fill (1,1) circle (1.6pt);
  \draw[<->,thick] (1,1.16) -- (1,3.84);
  \node[right=2pt] at (1,2.5) {$\Delta(D)=\tfrac12 r(D)$};
  \fill (5,5) circle (1.6pt);
  \node[below right=1pt] at (5,5) {\footnotesize $\bigl(\tfrac12,\tfrac12\bigr)$};
\end{tikzpicture}}
\caption{The achievable $(\rho,g)$ region at $q=0.2$. Record-independent
policies lie on the diagonal. The delayed threshold policy reaches the kink and
the maximum matched-mismatch increase $\Delta(D)$.}
\label{fig:theory}
\end{figure}

\begin{corollary}[Decay with observation-to-effect lag]\label{cor:imp}
For the symmetric Markov process,
$\Delta(D)=\tfrac12 e^{-D/\tau_c}$. Thus the increase is below a tolerance
$\eta$ once $D\ge\tau_c\ln(1/(2\eta))$. For a general process, zero
autocorrelation does not imply unpredictability from the full history. Appendix
\ref{app:beta} therefore gives a separate scalarized value bound using the
$\beta$-mixing coefficient.
\end{corollary}

\noindent\textbf{Provenance and scope of the exact result.}
The two-cell likelihood ordering is classical in opportunistic access
\cite{myopicaccess}, and the Markov mixing calculation is standard
\cite{lpw}. The contribution of Theorem~\ref{thm:limit} is the explicit
matched-mismatch achievable region and its vertical comparison with the best
record-independent policy. The equality is specific to the binary Markov model.
A non-Markov process can retain predictive information in older samples even
when $r(D)=0$. Appendix~\ref{app:beta} does not claim the same constrained
frontier for arbitrary capacity. It bounds a scalarized bounded objective by a
mixing coefficient.

The live square and memoryless profiles in Section~\ref{sec:eval} are qualitative
high- and low-predictability cases. Only the discrete-time Markov sweep in
Section~\ref{sec:eval-law} tests the exact curve in \eqref{eq:adv}.

\section{A Branch-Adaptive Controller}
\label{sec:controller}

The exact threshold policy needs the lag-$D$ predictive relation. A practical
controller can estimate the Markov flip probability from its delayed sample
stream and decide whether to run a feedback branch or a static branch. We use
AIMD as the feedback branch because it is simple and common. We do not claim that
AIMD attains Theorem~\ref{thm:limit}.

At decision epoch $t$, let $y_t$ be the newest binary capacity sample. It was
sampled $D_{\mathrm{obs}}$ earlier. Fixed delay does not change adjacent
transition counts, so the indicators $\mathbf 1[y_t\neq y_{t-1}]$ estimate the
one-step flip probability $p$. Let $c_t\in\{0,1,\bot\}$ be the newest resolved
outcome for a prior window, where $c_t=1$ denotes a capacity mismatch and
$c_t=0$ denotes no mismatch. The value $\bot$ means that no new outcome is
available. The controller keeps a real target $z_t\in[w_{\min},K]$. The executed
integer window is obtained by clipped unbiased rounding
$w_t=\mathcal R_K(z_t)$, where $\mathcal R_K(z)$ is either $\lfloor z\rfloor$ or
$\lceil z\rceil$ and has conditional mean $z$.

\begin{algorithm}[t]
\caption{Branch-adaptive admission for one contended unit}
\label{alg:adaptive}
\begin{algorithmic}[1]
\Require $K,w_{\min},\alpha,\beta,\theta,D,\bar a$
\State $z\gets w_{\min}$, $F\gets0$, $n\gets0$, $y_{\mathrm{prev}}\gets\bot$
\For{each decision epoch $t$}
  \State read newest sample $y_t$ and newest resolved outcome $c_t$
  \If{$y_{\mathrm{prev}}\neq\bot$}
    \State $F\gets F+\mathbf 1[y_t\neq y_{\mathrm{prev}}]$, $n\gets n+1$
  \EndIf
  \State $y_{\mathrm{prev}}\gets y_t$
  \State $\hat p\gets\min\{1/2,F/\max(n,1)\}$
  \State $\hat r\gets(1-2\hat p)^D$
  \If{$\hat r>\theta$}
    \If{$c_t=1$}
      \State $z\gets\max\{w_{\min},\beta z\}$
    \ElsIf{$c_t=0$}
      \State $z\gets\min\{z+\alpha,K\}$
    \EndIf
  \Else
    \State $z\gets\min\{\max\{\bar a,w_{\min}\},K\}$
  \EndIf
  \State $w_t\gets\mathcal R_K(z)$ and admit at most $w_t$ concurrent Jobs
\EndFor
\end{algorithmic}
\end{algorithm}

The cumulative estimator in Algorithm~\ref{alg:adaptive} supports the asymptotic
result below. A fixed sliding window is preferable under drift, but it leaves a
nonzero steady-state branch-error probability. The live experiment uses a
sliding window and is therefore a mechanical test rather than an evaluation of
the asymptotic rate.

\begin{proposition}[Continuous-target AIMD mismatch]\label{prop:bound}
Assume constant capacity $a$, a resolved mismatch signal exactly $D$ epochs after
an overshoot, a continuously valued target, and a backoff satisfying
$\beta(a+D\alpha)\le a$. While the feedback branch remains active, the target
forms a periodic sawtooth with mismatch fraction
\begin{equation}\label{eq:saw}
\rho_{\mathrm{aimd}}
=\frac{D\alpha}{(1-\beta)(a+D\alpha)}.
\end{equation}
Thus $\alpha\le\epsilon(1-\beta)a/[D(1-\epsilon(1-\beta))]$ is sufficient for
$\rho_{\mathrm{aimd}}\le\epsilon$ in this approximation.
\end{proposition}
\begin{proof}
The target rises from $\beta(a+D\alpha)$ to $a+D\alpha$ in a cycle of length
$(1-\beta)(a+D\alpha)/\alpha$. It exceeds $a$ for $D$ epochs. Their ratio gives
\eqref{eq:saw}. Solving the inequality for $\alpha$ gives the stated condition.
\end{proof}
Randomized rounding and variable outcome delay perturb this idealized cycle. The
proposition is an interpretation aid, not a safety guarantee for the live
implementation.

We next isolate branch selection from branch quality. Let $V_F(r)$ and $V_S(r)$
be the stationary per-epoch values of the supplied feedback and static branches
when the lag-$D$ correlation is $r$. Define
$d(r)=V_F(r)-V_S(r)$. Assume that in a neighborhood of a chosen switch level
$\theta\in(0,1)$, $d$ is differentiable, $d(\theta)=0$, and
$L=|d'(\theta)|>0$. This crossing assumption is part of the theorem. It is not
derived from the AIMD sawtooth proposition.

Let $p_\theta=(1-\theta^{1/D})/2$,
$s_\theta=2D\theta^{(D-1)/D}$, and
$\sigma_\theta^2=p_\theta(1-p_\theta)$. The comparator is
$V^{\mathrm{br}}(r)=\max\{V_F(r),V_S(r)\}$.

\begin{theorem}[Local branch-selection rate]\label{thm:adapt}
Consider the symmetric two-state Markov family in a shrinking neighborhood of
$p_\theta$, and let a selector observe the cumulative flip indicators and choose
between the two supplied branches. Over horizon $N$, the local minimax average
shortfall from $V^{\mathrm{br}}$ is
\begin{equation}\label{eq:rate}
\Theta\!\left(\frac{L s_\theta\sigma_\theta}{\sqrt N}\right),
\qquad
s_\theta\sigma_\theta
=D\theta^{(D-1)/D}\sqrt{1-\theta^{2/D}}.
\end{equation}
The branch selector in Algorithm~\ref{alg:adaptive} attains the upper rate with
cumulative counters. For fixed $\theta\in(0,1)$ and $D\to\infty$,
\begin{equation}\label{eq:largeD}
s_\theta\sigma_\theta
\sim \theta\sqrt{2D\log(1/\theta)}.
\end{equation}
For $D\ge2$, the sensitivity term is maximized over $\theta$ at
$\theta=((D-1)/D)^{D/2}$, with value
\begin{equation}\label{eq:sup}
\sqrt D\left(\frac{D-1}{D}\right)^{(D-1)/2}
\sim e^{-1/2}\sqrt D.
\end{equation}
For a fixed process with $r(D)\neq\theta$, the average branch-selection
shortfall is $O(1/N)$.
\end{theorem}
\begin{proof}[Proof sketch]
The delayed flip indicators are independent Bernoulli$(p)$ variables. Binomial
Bernstein concentration gives
$|\hat p_n-p|=O(\sigma_\theta/\sqrt n)$ locally. The derivative magnitude of
$p\mapsto(1-2p)^D$ at $p_\theta$ is $s_\theta$, so the correlation estimation
error is $O(s_\theta\sigma_\theta/\sqrt n)$. Local linearity of $d$ converts this
to branch loss, and averaging over the growing sample size yields the upper
bound. For the lower bound, two Bernoulli flip laws separated by
$\delta=\Theta(\sigma_\theta/\sqrt N)$ have bounded total Kullback-Leibler
divergence over $N$ samples. Le Cam's two-point argument \cite{tsybakov} then
forces loss of order $L s_\theta\sigma_\theta/\sqrt N$ on one of the two laws.
The exact sensitivity follows by substituting
$\sigma_\theta=\tfrac12\sqrt{1-\theta^{2/D}}$. Equation \eqref{eq:largeD} follows
from $\theta^{1/D}=1-\log(1/\theta)/D+o(D^{-1})$. Maximizing
$D u^{D-1}\sqrt{1-u^2}$ over $u=\theta^{1/D}$ gives
$u^2=(D-1)/D$ and \eqref{eq:sup}. Away from the crossing, the branch error
probability decays exponentially in the cumulative sample size, so total
shortfall is bounded and its average is $O(1/N)$.
\end{proof}

Theorem~\ref{thm:adapt} is a statistical decision result about choosing between
two branches. It does not show that $V_F$ equals the optimal frontier in
Theorem~\ref{thm:limit}. The experiments report both quantities separately.

\noindent\textbf{Scope.}
The exact frontier requires a symmetric binary Markov unit, stationarity and a
fixed observation-to-effect lag. Under a non-Markov process, older samples can
remain useful. Under drift, a sliding estimator trades tracking for a persistent
error floor. Random readable record ages require a cell-allocation problem rather
than simple averaging. Integer capacity also requires nested marginal decisions
to correspond to one scalar window. We leave those extensions open.
